\documentclass[11pt]{article}
\usepackage[margin=1in]{geometry}
\usepackage{amsmath, amssymb, amsthm}
\usepackage{hyperref}
\usepackage{enumitem}

\title{Math - The Superposition Principle and Undetermined Coefficients Revisited}
\author{Curated Problem Sets}
\date{}

\begin{document}

\maketitle

\tableofcontents
\newpage

\section{Concept Checklist \& Question Mapping}
To ensure comprehensive coverage of the chapter, the following concepts and problem types have been identified from the source material. Each problem in the set is mapped to one or more of these core concepts.

\begin{itemize}
    \item \textbf{C1: Superposition Application} - Scaling and adding known particular solutions without solving the differential equation from scratch. (Questions: 1, 2)
    \item \textbf{C2: Applicability Checks} - Identifying when the Method of Undetermined Coefficients can or cannot be used (e.g., fails for $1/t$, $\tan t$, $\ln t$). (Questions: 3, 4, 5)
    \item \textbf{C3: Full General Solution} - The standard process: finding $y_h$, guessing $y_p$, differentiating, substituting, and solving for coefficients. (Questions: 6, 7, 8, 9, 30)
    \item \textbf{C4: Initial Value Problems (IVP)} - Applying initial conditions to the \textit{full} general solution $y(t) = y_h(t) + y_p(t)$, not just $y_h(t)$. (Questions: 10, 11, 12)
    \item \textbf{C5: Form Determination \& Resonance ($s=1$)} - Guessing the form of $y_p$ and multiplying by $t$ when the forcing function's root matches a single root of the characteristic equation. (Questions: 13, 14, 15)
    \item \textbf{C6: Higher Multiplicity of Roots ($s=2$)} - Multiplying the guess by $t^2$ when the forcing function's root matches a double root of the characteristic equation. (Questions: 16, 17)
    \item \textbf{C7: The $a^t$ Trick} - Rewriting exponential functions of different bases (e.g., $10^t$) as $e^{t \ln a}$ to apply the method. (Questions: 18, 19)
    \item \textbf{C8: Missing Terms Trap} - Ensuring both sine and cosine terms are included in the guess even if only one appears in the forcing function. (Questions: 20, 21)
    \item \textbf{C9: "Close Enough" Root Trap} - Recognizing that roots must match exactly (e.g., $\pm 2i \neq \pm 3i$) for resonance to occur. (Questions: 22, 23)
    \item \textbf{C10: Reverse Engineering} - Reconstructing the homogeneous equation, characteristic polynomial, and forcing function given a set of solutions. (Questions: 24, 25)
    \item \textbf{C11: Real-World Modeling} - Applying the math to mass-spring systems or circuits, distinguishing between transient states ($y_h$) and steady-state responses ($y_p$). (Questions: 26, 27)
    \item \textbf{C12: Find the Flaw} - Diagnostic testing to identify subtle algebraic or conceptual errors in proposed solutions. (Questions: 28, 29)
\end{itemize}

\newpage

\section{Problem Sets}

\subsection*{Part A: Superposition and Applicability}
\begin{enumerate}
    \item Given that $y_1(t) = e^{-t}$ is a solution to $y'' + 2y' + y = 0$ (Wait, $e^{-t}$ solves the homogeneous, let's use a nonhomogeneous example). Let $L[y] = y'' + 3y' + 2y$. Given that $y_1(t) = t$ is a solution to $L[y] = 3t + 1$ and $y_2(t) = e^{3t}$ is a solution to $L[y] = 20e^{3t}$, use the superposition principle to find a particular solution to $y'' + 3y' + 2y = 9t + 3 - 10e^{3t}$.
    \item Let $L[y] = y'' - y$. Suppose $y_1(t) = -\frac{1}{2}\sin t$ is a solution to $L[y] = \sin t$, and $y_2(t) = \frac{1}{3}e^{2t}$ is a solution to $L[y] = e^{2t}$. Find a particular solution to $y'' - y = 4\sin t + 6e^{2t}$.
    \item Determine whether the method of undetermined coefficients can be applied to find a particular solution for $y'' - 4y = t^3 e^{2t} + \sin(5t)$. Do not solve.
    \item Determine whether the method of undetermined coefficients can be applied to find a particular solution for $y'' + 4y = \sec(2t) + t^2$. Do not solve.
    \item Determine whether the method of undetermined coefficients can be applied to find a particular solution for $y'' - 2y' + y = \frac{e^t}{t^2}$. Do not solve.
\end{enumerate}

\subsection*{Part B: Full General Solutions}
\begin{enumerate}[resume]
    \item Find the general solution to the differential equation: $y'' - 5y' + 6y = 2e^{4t}$.
    \item Find the general solution to the differential equation: $y'' + 3y' + 2y = 10\sin t$.
    \item Find the general solution to the differential equation: $y'' - 4y = 8t^2$.
    \item Find the general solution to the differential equation: $y'' + y' - 2y = e^t$.
    \item Find the general solution to the differential equation: $y'' - 2y' + 5y = 10$.
\end{enumerate}

\subsection*{Part C: Initial Value Problems}
\begin{enumerate}[resume]
    \item Solve the initial value problem: $y'' + 4y = 8$; $y(0) = 3$, $y'(0) = 4$.
    \item Solve the initial value problem: $y'' - y' = 2t$; $y(0) = 1$, $y'(0) = 0$.
    \item Solve the initial value problem: $y'' + y = \cos(2t)$; $y(0) = 0$, $y'(0) = 0$.
\end{enumerate}

\subsection*{Part D: Form Determination and Resonance}
\begin{enumerate}[resume]
    \item Determine the form of a particular solution $y_p$ for $y'' - 6y' + 9y = e^{2t} + \sin(3t)$. Do not solve for the coefficients.
    \item Determine the form of a particular solution $y_p$ for $y'' + 16y = t\cos(4t)$. Do not solve for the coefficients.
    \item Determine the form of a particular solution $y_p$ for $y'' - 5y' = t^2 + e^{5t}$. Do not solve for the coefficients.
    \item Determine the form of a particular solution $y_p$ for $y'' + 4y' + 4y = t e^{-2t}$. Do not solve for the coefficients.
    \item Determine the form of a particular solution $y_p$ for $y'' = 5t^2 - 3$. Do not solve for the coefficients.
\end{enumerate}

\subsection*{Part E: Advanced Tricks and Traps}
\begin{enumerate}[resume]
    \item Determine the form of a particular solution $y_p$ for $y'' - 3y' + 2y = 5^t$. Do not solve for the coefficients.
    \item Find the general solution to the differential equation: $y'' - y = 2^t$.
    \item Find the general solution to the differential equation: $y'' - 2y' + y = 3\cos t$.
    \item \textbf{Find the Flaw:} A student is asked to find the form of $y_p$ for $y'' + 9y = \cos(2t)$. The student writes: "The homogeneous roots are $r = \pm 3i$. The forcing function has $\cos(2t)$. Because cosine is in both the homogeneous solution and the forcing function, we have resonance. Guess $y_p = t(A\cos 2t + B\sin 2t)$." Identify the flaw and provide the correct form.
    \item Determine the form of a particular solution $y_p$ for $y'' - 4y' + 5y = e^{2t}\cos(3t)$. Do not solve for the coefficients.
\end{enumerate}

\subsection*{Part F: Reverse Engineering and Modeling}
\begin{enumerate}[resume]
    \item All that is known concerning a second-order constant-coefficient differential equation $y'' + py' + qy = g(t)$ is that $t + e^{-t}$, $t + e^{2t}$, and $t + e^{-t} + e^{2t}$ are solutions. Find the constants $p, q$ and the function $g(t)$.
    \item All that is known concerning a second-order constant-coefficient differential equation $y'' + py' + qy = g(t)$ is that $\sin(2t) + e^{-t}$, $\sin(2t) + te^{-t}$, and $\sin(2t) + e^{-t} + te^{-t}$ are solutions. Find the constants $p, q$ and the function $g(t)$.
    \item A mass-spring system is driven by an external force $F(t) = 10\cos t$. The mass is $m=1$, the damping coefficient is $c=3$, and the spring constant is $k=2$. The equation of motion is $y'' + 3y' + 2y = 10\cos t$. Find the steady-state solution (the particular solution) of the system.
    \item An LC circuit with no resistance (analogous to an undamped mass-spring system) is modeled by $q'' + 9q = 12\sin(3t)$, where $q(t)$ is the charge. Determine the form of the particular solution $q_p(t)$ and state the physical phenomenon that will occur in this circuit.
\end{enumerate}

\subsection*{Part G: Diagnostic Testing - "Find the Flaw"}
\begin{enumerate}[resume]
    \item \textbf{Find the Flaw:} Solve $y'' - y = t e^t$. \\
    \textit{Flawed Solution:} 1. $r^2 - 1 = 0 \implies r = \pm 1$. 2. Forcing is $t e^t$. The root is 1, which matches the homogeneous root, so we multiply by $t$. 3. Guess $y_p = A t e^t$. 4. $y_p' = A e^t + A t e^t$, $y_p'' = 2A e^t + A t e^t$. 5. Substitute: $(2A e^t + A t e^t) - A t e^t = t e^t \implies 2A e^t = t e^t$. This is impossible. \\
    \textbf{Your Task:} Find the flaw, explain it in one sentence, and provide the correct final solution.
    
    \item \textbf{Find the Flaw:} Solve $y'' + y = \sin(2t)$. \\
    \textit{Flawed Solution:} 1. $r^2 + 1 = 0 \implies r = \pm i$. 2. Forcing is $\sin(2t)$. 3. Guess $y_p = A\sin(2t)$. 4. $y_p' = 2A\cos(2t)$, $y_p'' = -4A\sin(2t)$. 5. Substitute: $-4A\sin(2t) + A\sin(2t) = \sin(2t) \implies -3A = 1 \implies A = -1/3$. 6. Final answer: $y = c_1\cos t + c_2\sin t - \frac{1}{3}\sin(2t)$. \\
    \textbf{Your Task:} Find the flaw, explain it in one sentence, and provide the correct final solution.
\end{enumerate}

\newpage

\section{Detailed Solutions}

\textbf{Solution 1:}
Let $L[y] = y'' + 3y' + 2y$. We are given $L[y_1] = 3t + 1$ and $L[y_2] = 20e^{3t}$. We want to solve $L[y_p] = 9t + 3 - 10e^{3t}$. Notice that $9t + 3 = 3(3t + 1)$ and $-10e^{3t} = -\frac{1}{2}(20e^{3t})$. By the Superposition Principle, $y_p(t) = 3y_1(t) - \frac{1}{2}y_2(t)$. Substituting the given functions: $y_p(t) = 3(t) - \frac{1}{2}(e^{3t}) = 3t - \frac{1}{2}e^{3t}$.

\textbf{Solution 2:}
We are given $L[y_1] = \sin t$ and $L[y_2] = e^{2t}$. We want to solve $L[y_p] = 4\sin t + 6e^{2t}$. By Superposition, $y_p(t) = 4y_1(t) + 6y_2(t) = 4(-\frac{1}{2}\sin t) + 6(\frac{1}{3}e^{2t}) = -2\sin t + 2e^{2t}$.

\textbf{Solution 3:}
Yes. The forcing function $f(t) = t^3 e^{2t} + \sin(5t)$ consists of a polynomial multiplied by an exponential, plus a sine function. These are all valid forms for the method of undetermined coefficients.

\textbf{Solution 4:}
No. The forcing function contains $\sec(2t) = \frac{1}{\cos(2t)}$. The method of undetermined coefficients cannot be applied to quotients of trigonometric functions because their derivatives do not form a finite set of linearly independent terms.

\textbf{Solution 5:}
No. The forcing function contains $t^{-2}$. Negative exponents are not polynomials, and their derivatives generate an infinite sequence of new terms ($t^{-3}, t^{-4}, \dots$).

\textbf{Solution 6:}
Step 1: Homogeneous solution. $r^2 - 5r + 6 = 0 \implies (r-2)(r-3) = 0 \implies r = 2, 3$. $y_h = c_1e^{2t} + c_2e^{3t}$. \\
Step 2: Particular solution. Guess $y_p = Ae^{4t}$. $y_p' = 4Ae^{4t}$, $y_p'' = 16Ae^{4t}$. \\
Step 3: Substitute. $16Ae^{4t} - 5(4Ae^{4t}) + 6(Ae^{4t}) = 2e^{4t} \implies (16 - 20 + 6)A = 2 \implies 2A = 2 \implies A = 1$. \\
Final Answer: $y(t) = c_1e^{2t} + c_2e^{3t} + e^{4t}$.

\textbf{Solution 7:}
Step 1: Homogeneous solution. $r^2 + 3r + 2 = 0 \implies (r+1)(r+2) = 0 \implies r = -1, -2$. $y_h = c_1e^{-t} + c_2e^{-2t}$. \\
Step 2: Particular solution. Guess $y_p = A\cos t + B\sin t$. $y_p' = -A\sin t + B\cos t$, $y_p'' = -A\cos t - B\sin t$. \\
Step 3: Substitute. $(-A\cos t - B\sin t) + 3(-A\sin t + B\cos t) + 2(A\cos t + B\sin t) = 10\sin t$. \\
Group terms: $\cos t(-A + 3B + 2A) + \sin t(-B - 3A + 2B) = 10\sin t \implies \cos t(A + 3B) + \sin t(-3A + B) = 10\sin t$. \\
Equate coefficients: $A + 3B = 0 \implies A = -3B$. $-3(-3B) + B = 10 \implies 10B = 10 \implies B = 1$. $A = -3$. \\
Final Answer: $y(t) = c_1e^{-t} + c_2e^{-2t} - 3\cos t + \sin t$.

\textbf{Solution 8:}
Step 1: Homogeneous solution. $r^2 - 4 = 0 \implies r = \pm 2$. $y_h = c_1e^{2t} + c_2e^{-2t}$. \\
Step 2: Particular solution. Guess $y_p = At^2 + Bt + C$. $y_p' = 2At + B$, $y_p'' = 2A$. \\
Step 3: Substitute. $2A - 4(At^2 + Bt + C) = 8t^2 \implies -4At^2 - 4Bt + (2A - 4C) = 8t^2$. \\
Equate coefficients: $-4A = 8 \implies A = -2$. $-4B = 0 \implies B = 0$. $2A - 4C = 0 \implies 2(-2) - 4C = 0 \implies -4 = 4C \implies C = -1$. \\
Final Answer: $y(t) = c_1e^{2t} + c_2e^{-2t} - 2t^2 - 1$.

\textbf{Solution 9:}
Step 1: Homogeneous solution. $r^2 + r - 2 = 0 \implies (r+2)(r-1) = 0 \implies r = 1, -2$. $y_h = c_1e^t + c_2e^{-2t}$. \\
Step 2: Particular solution. Forcing is $e^t$. Root is $\alpha = 1$. This matches a homogeneous root ($s=1$). Guess $y_p = Ate^t$. \\
Step 3: Differentiate. $y_p' = A(t+1)e^t$, $y_p'' = A(t+2)e^t$. \\
Step 4: Substitute. $A(t+2)e^t + A(t+1)e^t - 2Ate^t = e^t \implies e^t(At + 2A + At + A - 2At) = e^t \implies 3A = 1 \implies A = 1/3$. \\
Final Answer: $y(t) = c_1e^t + c_2e^{-2t} + \frac{1}{3}te^t$.

\textbf{Solution 10:}
Step 1: Homogeneous solution. $r^2 - 2r + 5 = 0 \implies r = \frac{2 \pm \sqrt{4 - 20}}{2} = 1 \pm 2i$. $y_h = c_1e^t\cos(2t) + c_2e^t\sin(2t)$. \\
Step 2: Particular solution. Guess $y_p = A$. $y_p' = 0, y_p'' = 0$. Substitute: $0 - 0 + 5A = 10 \implies A = 2$. \\
Final Answer: $y(t) = c_1e^t\cos(2t) + c_2e^t\sin(2t) + 2$.

\textbf{Solution 11:}
Step 1: General Solution. $r^2 + 4 = 0 \implies r = \pm 2i$. $y_h = c_1\cos(2t) + c_2\sin(2t)$. Guess $y_p = A$. $4A = 8 \implies A = 2$. $y(t) = c_1\cos(2t) + c_2\sin(2t) + 2$. \\
Step 2: Apply ICs. $y(0) = c_1 + 2 = 3 \implies c_1 = 1$. $y'(t) = -2c_1\sin(2t) + 2c_2\cos(2t)$. $y'(0) = 2c_2 = 4 \implies c_2 = 2$. \\
Final Answer: $y(t) = \cos(2t) + 2\sin(2t) + 2$.

\textbf{Solution 12:}
Step 1: General Solution. $r^2 - r = 0 \implies r(r-1) = 0 \implies r = 0, 1$. $y_h = c_1 + c_2e^t$. Forcing is $2t$ (polynomial degree 1, root 0). Since 0 is a root of the characteristic equation, $s=1$. Guess $y_p = t(At + B) = At^2 + Bt$. $y_p' = 2At + B$, $y_p'' = 2A$. Substitute: $2A - (2At + B) = 2t \implies -2At + (2A - B) = 2t \implies -2A = 2 \implies A = -1$. $2A - B = 0 \implies -2 - B = 0 \implies B = -2$. $y(t) = c_1 + c_2e^t - t^2 - 2t$. \\
Step 2: Apply ICs. $y(0) = c_1 + c_2 = 1$. $y'(t) = c_2e^t - 2t - 2$. $y'(0) = c_2 - 2 = 0 \implies c_2 = 2$. $c_1 + 2 = 1 \implies c_1 = -1$. \\
Final Answer: $y(t) = -1 + 2e^t - t^2 - 2t$.

\textbf{Solution 13:}
Step 1: General Solution. $r^2 + 1 = 0 \implies r = \pm i$. $y_h = c_1\cos t + c_2\sin t$. Forcing is $\cos(2t)$. Roots are $\pm 2i$. No resonance. Guess $y_p = A\cos(2t) + B\sin(2t)$. $y_p'' = -4A\cos(2t) - 4B\sin(2t)$. Substitute: $-3A\cos(2t) - 3B\sin(2t) = \cos(2t) \implies A = -1/3, B = 0$. $y(t) = c_1\cos t + c_2\sin t - \frac{1}{3}\cos(2t)$. \\
Step 2: Apply ICs. $y(0) = c_1 - 1/3 = 0 \implies c_1 = 1/3$. $y'(t) = -c_1\sin t + c_2\cos t + \frac{2}{3}\sin(2t)$. $y'(0) = c_2 = 0$. \\
Final Answer: $y(t) = \frac{1}{3}\cos t - \frac{1}{3}\cos(2t)$.

\textbf{Solution 14:}
Homogeneous roots: $r^2 - 6r + 9 = 0 \implies (r-3)^2 = 0 \implies r = 3$ (multiplicity 2). Forcing $f_1 = e^{2t}$ has root 2 (no match). Forcing $f_2 = \sin(3t)$ has roots $\pm 3i$ (no match). \\
Final Answer: $y_p(t) = Ae^{2t} + B\cos(3t) + C\sin(3t)$.

\textbf{Solution 15:}
Homogeneous roots: $r^2 + 16 = 0 \implies r = \pm 4i$. Forcing is $t\cos(4t)$, which has roots $\pm 4i$. This matches the homogeneous roots exactly once ($s=1$). The polynomial is degree 1. \\
Final Answer: $y_p(t) = t((At + B)\cos(4t) + (Ct + D)\sin(4t))$.

\textbf{Solution 16:}
Homogeneous roots: $r^2 - 5r = 0 \implies r = 0, 5$. Forcing $f_1 = t^2$ has root 0. Matches $r=0$ ($s=1$). Base guess $At^2+Bt+C \implies t(At^2+Bt+C)$. Forcing $f_2 = e^{5t}$ has root 5. Matches $r=5$ ($s=1$). Base guess $De^{5t} \implies Dte^{5t}$. \\
Final Answer: $y_p(t) = At^3 + Bt^2 + Ct + Dte^{5t}$.

\textbf{Solution 17:}
Homogeneous roots: $r^2 + 4r + 4 = 0 \implies (r+2)^2 = 0 \implies r = -2$ (multiplicity 2). Forcing is $te^{-2t}$, root is $-2$. Matches the double root, so $s=2$. Base guess is $(At+B)e^{-2t}$. \\
Final Answer: $y_p(t) = t^2(At+B)e^{-2t} = (At^3 + Bt^2)e^{-2t}$.

\textbf{Solution 18:}
Homogeneous roots: $r^2 = 0 \implies r = 0$ (multiplicity 2). Forcing is $5t^2 - 3$, root is 0. Matches the double root, so $s=2$. Base guess is $At^2+Bt+C$. \\
Final Answer: $y_p(t) = t^2(At^2+Bt+C) = At^4 + Bt^3 + Ct^2$.

\textbf{Solution 19:}
Rewrite $5^t = (e^{\ln 5})^t = e^{t\ln 5}$. Homogeneous roots: $r^2 - 3r + 2 = 0 \implies r = 1, 2$. The forcing root is $\ln 5 \approx 1.609$, which does not match 1 or 2. \\
Final Answer: $y_p(t) = Ae^{t\ln 5}$ or $A5^t$.

\textbf{Solution 20:}
Rewrite $2^t = e^{t\ln 2}$. Homogeneous roots: $r^2 - 1 = 0 \implies r = \pm 1$. Guess $y_p = Ae^{t\ln 2}$. $y_p'' = A(\ln 2)^2 e^{t\ln 2}$. Substitute: $A(\ln 2)^2 e^{t\ln 2} - Ae^{t\ln 2} = e^{t\ln 2} \implies A((\ln 2)^2 - 1) = 1 \implies A = \frac{1}{(\ln 2)^2 - 1}$. \\
Final Answer: $y(t) = c_1e^t + c_2e^{-t} + \frac{2^t}{(\ln 2)^2 - 1}$.

\textbf{Solution 21:}
Homogeneous roots: $r^2 - 2r + 1 = 0 \implies r = 1$ (mult 2). $y_h = c_1e^t + c_2te^t$. Guess $y_p = A\cos t + B\sin t$. $y_p' = -A\sin t + B\cos t$, $y_p'' = -A\cos t - B\sin t$. Substitute: $(-A\cos t - B\sin t) - 2(-A\sin t + B\cos t) + (A\cos t + B\sin t) = 3\cos t \implies -2B\cos t + 2A\sin t = 3\cos t \implies -2B = 3 \implies B = -3/2$, and $2A = 0 \implies A = 0$. \\
Final Answer: $y(t) = c_1e^t + c_2te^t - \frac{3}{2}\sin t$.

\textbf{Solution 22:}
\textbf{Flaw:} The student confused the frequency of the forcing function ($\beta = 2$) with the homogeneous roots ($r = \pm 3i$); since $\pm 2i \neq \pm 3i$, there is no resonance. \\
\textbf{Correct Form:} $y_p(t) = A\cos(2t) + B\sin(2t)$.

\textbf{Solution 23:}
Homogeneous roots: $r^2 - 4r + 5 = 0 \implies r = \frac{4 \pm \sqrt{16 - 20}}{2} = 2 \pm i$. Forcing function is $e^{2t}\cos(3t)$, which corresponds to roots $2 \pm 3i$. Since $2 \pm i \neq 2 \pm 3i$, there is no resonance ($s=0$). \\
Final Answer: $y_p(t) = e^{2t}(A\cos(3t) + B\sin(3t))$.

\textbf{Solution 24:}
Let $y_1 = t + e^{-t}$, $y_2 = t + e^{2t}$, $y_3 = t + e^{-t} + e^{2t}$. The difference between any two solutions is a homogeneous solution. $y_2 - y_1 = e^{2t} - e^{-t}$. $y_3 - y_2 = e^{-t}$. $y_3 - y_1 = e^{2t}$. The fundamental set of homogeneous solutions is $\{e^{-t}, e^{2t}\}$. The roots are $r = -1, 2$. The characteristic polynomial is $(r+1)(r-2) = r^2 - r - 2 = 0$. Thus, $p = -1, q = -2$. The remaining part of the solutions is $y_p = t$. Plug $y_p$ into the left side: $y_p'' - y_p' - 2y_p = 0 - 1 - 2(t) = -2t - 1$. \\
Final Answer: $p = -1, q = -2, g(t) = -2t - 1$.

\textbf{Solution 25:}
Let $y_1 = \sin(2t) + e^{-t}$, $y_2 = \sin(2t) + te^{-t}$, $y_3 = \sin(2t) + e^{-t} + te^{-t}$. Differences yield homogeneous solutions $e^{-t}$ and $te^{-t}$. This implies a repeated root $r = -1$. The characteristic polynomial is $(r+1)^2 = r^2 + 2r + 1 = 0$. Thus, $p = 2, q = 1$. The particular solution is $y_p = \sin(2t)$. Plug $y_p$ into the DE: $y_p'' + 2y_p' + y_p = -4\sin(2t) + 2(2\cos(2t)) + \sin(2t) = -3\sin(2t) + 4\cos(2t)$. \\
Final Answer: $p = 2, q = 1, g(t) = -3\sin(2t) + 4\cos(2t)$.

\textbf{Solution 26:}
The steady-state solution is the particular solution $y_p$. Guess $y_p = A\cos t + B\sin t$. $y_p' = -A\sin t + B\cos t$, $y_p'' = -A\cos t - B\sin t$. Substitute into $y'' + 3y' + 2y = 10\cos t$: $(-A\cos t - B\sin t) + 3(-A\sin t + B\cos t) + 2(A\cos t + B\sin t) = 10\cos t$. Grouping yields $(A + 3B)\cos t + (-3A + B)\sin t = 10\cos t$. System: $A + 3B = 10$ and $-3A + B = 0 \implies B = 3A$. Substitute: $A + 3(3A) = 10 \implies 10A = 10 \implies A = 1$. $B = 3(1) = 3$. \\
Final Answer: $y_p(t) = \cos t + 3\sin t$.

\textbf{Solution 27:}
Homogeneous roots: $r^2 + 9 = 0 \implies r = \pm 3i$. The forcing function is $12\sin(3t)$, which has roots $\pm 3i$. Because the forcing frequency matches the natural frequency of the undamped system, we have resonance ($s=1$). The form is $q_p(t) = t(A\cos(3t) + B\sin(3t))$. \\
Final Answer: Form is $q_p(t) = t(A\cos(3t) + B\sin(3t))$. The physical phenomenon is \textbf{pure resonance}, where the amplitude of the charge oscillates and grows unboundedly over time.

\textbf{Solution 28:}
\textbf{Flaw:} The student forgot that the polynomial part of the forcing function $t e^t$ is degree 1, so the base guess must be $(At+B)e^t$, which becomes $(At^2+Bt)e^t$ after multiplying by $t$ for resonance. \\
\textbf{Correct Solution:} Guess $y_p = (At^2 + Bt)e^t$. $y_p' = (2At + B)e^t + (At^2 + Bt)e^t = (At^2 + (2A+B)t + B)e^t$. $y_p'' = (2At + 2A + B)e^t + (At^2 + (2A+B)t + B)e^t = (At^2 + (4A+B)t + 2A + 2B)e^t$. Substitute: $y_p'' - y_p = (4At + 2A + 2B)e^t = te^t$. Equate coefficients: $4A = 1 \implies A = 1/4$. $2A + 2B = 0 \implies 2(1/4) + 2B = 0 \implies B = -1/4$. Final answer: $y_p(t) = (\frac{1}{4}t^2 - \frac{1}{4}t)e^t$.

\textbf{Solution 29:}
\textbf{Flaw:} The student assumed that because the forcing function is $\sin(2t)$, the guess only needs to contain $\sin(2t)$, forgetting the "Missing Terms" rule that both sine and cosine must be included. \\
\textbf{Correct Solution:} Guess $y_p = A\cos(2t) + B\sin(2t)$. $y_p'' = -4A\cos(2t) - 4B\sin(2t)$. Substitute: $(-4A\cos(2t) - 4B\sin(2t)) + (A\cos(2t) + B\sin(2t)) = \sin(2t) \implies -3A\cos(2t) - 3B\sin(2t) = \sin(2t)$. Equate coefficients: $-3A = 0 \implies A = 0$. $-3B = 1 \implies B = -1/3$. Final answer: $y(t) = c_1\cos t + c_2\sin t - \frac{1}{3}\sin(2t)$.

\newpage

\section{Concept Checklist \& Question Mapping (Reference)}
\begin{itemize}
    \item \textbf{C1: Superposition Application} - Questions: 1, 2
    \item \textbf{C2: Applicability Checks} - Questions: 3, 4, 5
    \item \textbf{C3: Full General Solution} - Questions: 6, 7, 8, 9, 30
    \item \textbf{C4: Initial Value Problems (IVP)} - Questions: 10, 11, 12
    \item \textbf{C5: Form Determination \& Resonance ($s=1$)} - Questions: 13, 14, 15
    \item \textbf{C6: Higher Multiplicity of Roots ($s=2$)} - Questions: 16, 17
    \item \textbf{C7: The $a^t$ Trick} - Questions: 18, 19
    \item \textbf{C8: Missing Terms Trap} - Questions: 20, 21
    \item \textbf{C9: "Close Enough" Root Trap} - Questions: 22, 23
    \item \textbf{C10: Reverse Engineering} - Questions: 24, 25
    \item \textbf{C11: Real-World Modeling} - Questions: 26, 27
    \item \textbf{C12: Find the Flaw} - Questions: 28, 29
\end{itemize}

\end{document}